\section{notes for conclusion}

``Hayek's central argument... was not so much that a socialist economy could not transmit the necessary data, but rather that it could not generate it to begin with. Without the competitive process of discovery and innovation, a socialist economy would have available only a small fraction of the knowledge that is utilized in a competitive economy. The task faced by proponents of market socialism is to explain exactly how spontaneous discovery is to occur within a planned economic system.'' (Chamberlain 1998)

Chamberlain, A., 1998. The Socialist Economic Calculation Debate and the Austrian Critique of Central Planning. The Idea Shop. Available at: http://the-idea-shop.com/article/97/the-socialist-economic-calculation-debate-and-the-austrian-critique-of-central-planning [Accessed April 9, 2008].

Tie to point about metis: it may simply be that computer agents can never capture the core of the processes intrinsic to real human beings setting price. Two consequences from that: high-level economic arguments may, paradoxically, capture the aggregate behaviour of the system very well - not because they are built on micro-foundations, but because people's complex economic choices, in some circumstances, \emph{do} aggregate.

The second would be that \gls{abm} will struggle unless it accepts its limitations and stops placing psychological credibility high on its list of priorities.